\section{Introduction}
\label{sec:intro}

Humans come to know an unfamiliar physical world by acting in it. They push, drop, and collide objects, form an understanding of the principles behind what they see, and then carry that understanding back to vision: a glance at a new scene is enough to picture how it will unfold. We ask whether multimodal agents can learn the physical world through interaction in the same way. Improving on a practiced task is not enough to show this. An agent can improve over a few trials by searching a small action space, replaying a trajectory that worked, or exploiting a visual cue that happens to correlate with success. Learning in the sense we study means that what the agent gathers by interacting with a world becomes an understanding of that world, one that still holds when the agent looks at a new scene from it.

Current benchmarks leave this distinction partly hidden. Passive physical-reasoning benchmarks ask models to judge physical concepts or outcomes from fixed images, videos, or questions; they do not test what feedback from action changes. Interactive physics benchmarks such as PHYRE~\cite{phyre} and I-PHYRE~\cite{iphyre}, and their recent VLM-agent counterparts~\cite{deepphy,ipr1}, measure task success, but success on the acted-on instance can mix physical inference with search. Indeed, VLMs trained through interaction improve on their training task without acquiring physical intuitions that transfer to related tasks~\cite{schulzebuschoff2026}. Scientific-discovery benchmarks~\cite{physgym,newtonbench,discoverphysics} let agents probe worlds with hidden or altered laws, but their observations and outputs are numeric or symbolic. Visual benchmarks of hidden physics~\cite{comphy,inphyre} test transfer to new scenes, but supply the evidence instead of letting the agent collect it. What remains open is the full path from interaction to vision: whether an agent that has experimented in a physical world can foresee how a new scene of that world will look.

Our benchmark is built on one observation: interacting with a physical world and seeing it correctly are two sides of the same visual experiment. An experiment pairs a rendered scene, an intervention, and a time interval with the rendered outcome. During \emph{interaction}, the agent chooses experiments and the world shows their outcomes. During \emph{evaluation}, the roles reverse: the world poses a new experiment on an unseen scene, and the agent must foresee its outcome by choosing the correct future frame among four rendered candidates. Because the two phases share the same scenes, interventions, and renderer, the evaluation asks the agent to use what it learned in exactly the form in which it learned it.

Each latent physical world is indexed by four factors: a physical-phenomenon family, a world-parameter regime of shared constants, a latent-law condition, and a latent-law realization. The condition specifies how the world departs from canonical physics, through a structural change to the governing law, a hidden coefficient field, or both, so prior knowledge of textbook physics is not enough. Each realization owns fifteen scenes, split into exploration scenes for the agent's own experiments, checkpoint scenes that track learning during interaction, and held-out scenes for the final evaluation. During interaction, the agent spends a fixed budget of experiments; at fixed points, it predicts the outcome of a host-posed checkpoint experiment before seeing it. After interaction, all tools and feedback are frozen, and the agent answers each held-out experiment from its interaction history alone.

Learning is then measured as what interaction changes. The checkpoint predictions trace an online learning curve within each episode. On the held-out evaluation, a budget sweep compares agents with and without interaction, and a cross-realization control replaces the agent's own interaction history with one from a sibling world that shares the family, regime, and condition but follows a different hidden law. If accuracy does not drop, the history was not used as evidence about this particular world. A canonical-physics baseline identifies held-out experiments that textbook physics alone can answer. We organize the analysis around three questions. Does interaction improve how the agent sees new scenes (RQ1)? Is the improvement specific to the world's own laws, and how does it differ between structural and coefficient-field changes (RQ2)? Where does learning break: in gathering evidence, forming a law, or applying it to images (RQ3)? For RQ3, a measurement probe asks the agent to estimate a hidden intrinsic attribute, such as mass or friction, with basic experiment tools, and law-informed ceilings measure how well the agent sees new scenes once it is told the law.

Our contributions are:
\begin{itemize}
  \item A benchmark that frames learning the physical world through interaction as a visual experiment seen from two sides: the agent experiments and observes during interaction, then foresees rendered outcomes of new experiments with interaction frozen.
  \item A five-level hierarchy of latent physical worlds that separates phenomenon families, shared constants, structural law changes, hidden coefficient fields, and individual scenes, with simulator-generated ground truth and rendered candidates. \TODO{State the number of families and instances validated end to end; non-central families currently ignore the hidden law in their ground truth.}
  \item An evaluation protocol that measures learning as what interaction changes, through online learning curves, a support-budget sweep, and same- versus cross-realization evaluation, and that locates failures with a measurement probe, canonical-physics baselines, and law-informed ceilings. We apply it to frontier multimodal agents at the model--harness level. \TODO{Replace this clause with the main empirical finding once results land.}
\end{itemize}

\begin{figure*}[t]
  \centering
  % \includegraphics[width=\textwidth]{figures/fig1_hero.pdf}
  \fbox{\parbox[c][6cm][c]{0.95\textwidth}{\centering
    \textbf{Figure 1 placeholder (hero figure).}\\[0.5em]
    \emph{(a) World:} the five-level hierarchy, with one path expanded (e.g., a projectile world with a hidden gravity field).\\
    \emph{(b) Interact:} the agent chooses experiments and sees rendered outcomes; checkpoint predictions trace a learning curve.\\
    \emph{(c) Foresee:} interaction is frozen; on a new scene, the agent chooses the correct future frame among four candidates, one of which follows textbook physics.
  }}
  \caption{\textbf{Learning a physical world through interaction.} Interaction and evaluation are the same visual experiment seen from two sides. During interaction, the agent poses experiments and the world shows their outcomes. During evaluation, the world poses new experiments on unseen scenes, and the agent must foresee their outcomes from what it learned.}
  \label{fig:hero}
\end{figure*}

Section~\ref{sec:related} positions the benchmark relative to prior work. Section~\ref{sec:benchmark} defines the visual experiment, the world hierarchy, and the episode. Section~\ref{sec:construction} describes the construction of worlds, tools, and renders. Section~\ref{sec:protocol} defines the metrics and controls. Section~\ref{sec:validation} gives validation checks, Section~\ref{sec:experiments} describes the experiments and analyses, and Section~\ref{sec:conclusion} discusses limitations and conclusions.
